% --- Archon physics formalization source begin ---
% archon:physics
% archon:chemistry
% archon:covers IChO2026Problems/problem_icho_2026_t9_a1.lean
% archon:source-report reports/icho_2026/problem_icho_2026_t9_a1.source.json
% archon:problem-id icho_2026_t9
% archon:part-id T9-A1

\chapter{Icho Chemistry problem icho\_2026\_t9\_a1}
\label{ch:IChO2026Problems_problem_icho_2026_t9_a1}

\paragraph{Problem source.}
\#\# Source
58th International Chemistry Olympiad, Tashkent, Uzbekistan, 2026.
Official-materials index: https://www.icho2026.uz/problems
English problem PDF: ../icho\_2026\_source/raw/theory\_problem.pdf
English solution/rubric PDF: ../icho\_2026\_source/raw/theory\_solution.pdf
Problem PDF page 84; printed local question page 1; solution starts on PDF page 80.
The page PNG is primary visual evidence for formulas, structures, tables, and figures.

\#\# T9. Cyclodextrin Chemistry
T9. Cyclodextrin Chemistry 7\% Cyclodextrins (CD) are a family of cyclic oligosaccharides, consisting of glucose subunits joined by α1,4-glycosidic bonds. The three most common cyclodextrins (α-, β-, γ-cyclodextrin) contain 6, 7, and 8 α-D-glucopyranoside units, respectively.

\#\# Current subquestion 9.1
Calculate the molar mass of β-CD. Assume M𝑤(glucose) = 180.16 g mol−1.

\#\# Dataset metadata
IChO 2026; T9; raw rubric points=2; kind=theory; formalization\_ready=true.

\paragraph{Shared problem context.}
T9. Cyclodextrin Chemistry 7\% Cyclodextrins (CD) are a family of cyclic oligosaccharides, consisting of glucose subunits joined by α1,4-glycosidic bonds. The three most common cyclodextrins (α-, β-, γ-cyclodextrin) contain 6, 7, and 8 α-D-glucopyranoside units, respectively.

\paragraph{Current subquestion.}
Calculate the molar mass of β-CD. Assume M𝑤(glucose) = 180.16 g mol−1.

\paragraph{Recorded answer/context.}
𝑀𝑤(𝛽−𝐶𝑦𝑐𝑙𝑜𝑑𝑒𝑥𝑡𝑟𝑖𝑛) = 180.16 ⋅7–18.016 ⋅7 = 1135.01g/mol Mw: β-CD = 1135.01 g/mol 2 points 2.0 pt

\paragraph{Figure/image path.}
../icho\_2026\_source/image/T9\_page-1.png

\paragraph{Formalization target.}
The verified Lean formalization is in `IChO2026Problems/problem\_icho\_2026\_t9\_a1.lean`,
with its theorem and lemma proofs complete.
The Lean declarations must preserve every mathematical object, quantifier, hypothesis, side condition, approximation/error guarantee, and requested conclusion from the source statement.
The implementation reuses the pinned Mathlib, Physlib, and Chemistry libraries,
with local definitions only for source-specific objects.

\begin{theorem}[IChO 2026 T9-A1: molar mass of $\beta$-cyclodextrin]
\label{thm:physics:icho_2026_t9_a1:target}
\lean{IChO2026.T9.A1.betaCD_molarMass_exact, IChO2026.T9.A1.betaCD_molarMass_reported}
\leanok
The molar mass of $\beta$-cyclodextrin, computed from the condensation
stoichiometry of the cyclic $\alpha$-1,4-glucan
($n = 7$ glucose units, one water eliminated per glycosidic bond, $n$ bonds
in the ring), is exactly
$M_w(\beta\text{-CD}) = 7 \cdot 180.16 - 7 \cdot 18.016 = 1135.008~\mathrm{g\,mol^{-1}}$,
which rounds to the reported $1135.01~\mathrm{g\,mol^{-1}}$ of the official
marking scheme.
\end{theorem}
\begin{proof}
\leanok
Unfold the local stoichiometric interface
(\texttt{Cyclodextrin.molarMass}, \texttt{cyclicGlucanMolarMass},
\texttt{cyclicGlycosidicBonds}, \texttt{glucoseMolarMass},
\texttt{waterMolarMass}, \texttt{Cyclodextrin.glucoseUnits}) and close the
resulting numeral identity over $\mathbb{R}$ by \texttt{norm\_num}; the
reported value follows by \texttt{round\_eq} and \texttt{Int.floor\_eq\_iff},
since $\lfloor 1135.008 \cdot 100 + 1/2 \rfloor = 113501$.
\end{proof}
% --- Archon physics formalization source end ---
